\documentclass[12pt,a4paper]{article}
\usepackage[utf8]{inputenc}
\usepackage[T1]{fontenc}
\usepackage[french]{babel}
\usepackage{amsmath,amssymb}
\usepackage{geometry}
\geometry{margin=2.5cm}
\usepackage{enumitem}
\usepackage{parskip}
\usepackage{array}
\usepackage{booktabs}
\usepackage{eurosym}   % pour \euro
\usepackage{xcolor} 

\begin{document}
	
\huge\textbf{TS2BIOAL} 

\large$\color{white}{A}\qquad$ \textbf{Évaluation formative : Évolutions}\normalsize 
	
\section*{Pénalités hors barème (jusqu'à \(-2\) points)}

\begin{center}
	\renewcommand{\arraystretch}{1.5}
	\begin{tabular}{|>{\raggedright\arraybackslash}p{4.2cm}|>{\raggedright\arraybackslash}p{9cm}|c|}
		\hline
		\textbf{Critère} & \textbf{Indicateurs} & \textbf{Malus max} \\
		\hline
		Rédaction & Absence de phrase réponse, pas de justification, calculs non explicités & \(-0{,}5\) \\
		\hline
		Orthographe & Fautes d'orthographe d'usage ou grammaticale & \(-0{,}5\) \\
		\hline
		Grammaire / Syntaxe & Phrases incorrectes, accords manquants, ponctuation absente & \(-0{,}5\) \\
		\hline
		Graphie & Écriture illisible, chiffres ambigus, ratures non assumées & \(-0{,}25\) \\
		\hline
		Présentation générale & Copie sale, exercices non identifiés, absence de numérotation & \(-0{,}25\) \\
		\hline
		\multicolumn{2}{|c|}{\textbf{Total pénalités}} & \(\mathbf{-2}\) \\
		\hline
	\end{tabular}
\end{center}

\vspace{0.3cm}

\begin{center}
	\fbox{%
		\parbox{0.9\linewidth}{%
			\textbf{Calcul de la note finale :}\\[0.2cm]
			Note finale \(=\) Total barème (sur 20) \(-\) Total pénalités.\\[0.2cm]
			Si le résultat est négatif, la note est ramenée à \(0\).%
		}%
	}
\end{center}

\vspace{0.5cm}
\section*{Barème pondéré -- Note finale sur 20}

\begin{center}
\renewcommand{\arraystretch}{1.5}
\begin{tabular}{lcc}
\hline
\textbf{Exercice} & \textbf{Coefficient} & \textbf{Points} \\
\hline
Exercice 1 -- Facile & 2 & 4 \\
\hline
Exercice 2 -- Moyen & 3 & 6 \\
\hline
Exercice 3 -- Difficile & 5 & 10 \\
\hline
\textbf{Total} & \textbf{10} & \textbf{20} \\
\hline
\end{tabular}
\end{center}

\section*{Exercice 1 -- Notions de base (facile) -- 4 points}

Un commerçant vend un article à \(80\)~\euro{}.

\begin{enumerate}[label=\textbf{\arabic*.}]
    \item Le prix augmente de \(15\%\). Calculer le coefficient multiplicateur associé, puis le nouveau prix.
    \item Le prix de cet article passe ensuite de \(92\)~\euro{} à \(78{,}20\)~\euro{}. Calculer le taux d'évolution en pourcentage. Interpréter le signe.
    \item Un autre article subit une baisse de \(8\%\). Son prix final est \(46\)~\euro{}. Retrouver son prix initial.
\end{enumerate}

\begin{center}
\renewcommand{\arraystretch}{1.5}
\begin{tabular}{>{\raggedright\arraybackslash}p{10.5cm} c}
\hline
\textbf{Compétence mathématique} & \textbf{Points} \\
\hline
Calculer un coefficient multiplicateur à partir d'un taux & 1 \\
\hline
Appliquer la relation \(V_f = CM \times V_i\) & 1 \\
\hline
Calculer un taux d'évolution à partir de \(V_i\) et \(V_f\) & 1 \\
\hline
Interpréter un taux d'évolution & 0{,}5 \\
\hline
Déterminer une valeur initiale connaissant \(V_f\) et \(CM\) & 0{,}5 \\
\hline
\textbf{Total Exercice 1} & \textbf{4} \\
\hline
\end{tabular}
\end{center}

\section*{Exercice 2 -- Évolutions successives et taux global (moyen) -- 6 points}

Une entreprise compte \(15\,000\) clients au 1\textsuperscript{er} janvier 2020.

\begin{itemize}
    \item En 2020, elle perd \(6\%\) de clients.
    \item En 2021, elle gagne \(10\%\).
    \item En 2022, elle perd \(4\%\).
\end{itemize}

\begin{enumerate}[label=\textbf{\arabic*.}]
    \item Donner les coefficients multiplicateurs \(CM_1\), \(CM_2\) et \(CM_3\).
    \item Calculer le coefficient multiplicateur global.
    \item En déduire le taux global d'évolution sur les trois années.
    \item Calculer le nombre de clients fin 2022 (arrondir à l'unité).
    \item Un stagiaire affirme : « Le taux global est \(-6+10-4 = 0\%\), donc le nombre de clients n'a pas changé. » Expliquer pourquoi il a tort.
\end{enumerate}

\begin{center}
\renewcommand{\arraystretch}{1.5}
\begin{tabular}{>{\raggedright\arraybackslash}p{10.5cm} c}
\hline
\textbf{Compétence mathématique} & \textbf{Points} \\
\hline
Convertir des taux en coefficients multiplicateurs & 1 \\
\hline
Multiplier des coefficients multiplicateurs pour obtenir un coefficient global & 1{,}5 \\
\hline
Calculer un taux global à partir du coefficient global & 1 \\
\hline
Calculer une valeur finale après évolutions successives & 1{,}5 \\
\hline
Argumenter sur la non-additivité des taux & 1 \\
\hline
\textbf{Total Exercice 2} & \textbf{6} \\
\hline
\end{tabular}
\end{center}

\section*{Exercice 3 -- Taux moyen et taux inconnu (difficile) -- 10 points}

Un capital de \(2\,000\)~\euro{} est placé. Il subit les évolutions annuelles suivantes :
\[
+10\%,\quad -5\%,\quad +8\%,\quad -3\%,\quad +6\%.
\]

\begin{enumerate}[label=\textbf{\arabic*.}]
    \item Calculer le coefficient multiplicateur global sur ces 5 années.
    \item Calculer le taux global d'évolution en pourcentage (arrondir à \(0{,}01\%\)).
    \item Calculer le taux annuel moyen équivalent (arrondir à \(0{,}01\%\)).
    \item Calculer le capital final (arrondir au centime).
    \item Un autre placement a augmenté de \(44\%\) en deux ans, avec le même taux annuel \(t\) chaque année. Déterminer \(t\).
\end{enumerate}

\begin{center}
\renewcommand{\arraystretch}{1.5}
\begin{tabular}{>{\raggedright\arraybackslash}p{10.5cm} c}
\hline
\textbf{Compétence mathématique} & \textbf{Points} \\
\hline
Calculer un coefficient multiplicateur global & 2 \\
\hline
Calculer un taux global en pourcentage & 1{,}5 \\
\hline
Calculer un taux moyen à partir du coefficient global et du nombre de périodes & 3 \\
\hline
Calculer une valeur finale à partir de \(V_i\) et \(CM_{\text{global}}\) & 1{,}5 \\
\hline
Résoudre une équation de la forme \(\left(1+\frac{t}{100}\right)^n = CM\) pour trouver un taux inconnu & 2 \\
\hline
\textbf{Total Exercice 3} & \textbf{10} \\
\hline
\end{tabular}
\end{center}

\vspace{0.5cm}
\begin{center}
\textbf{Note finale : 4 + 6 + 10 = 20 points}
\end{center}

% ============================================================
% CORRECTION DÉTAILLÉE
% ============================================================

\newpage
\section*{Correction détaillée}

\paragraph{Rappels généraux.}
\begin{itemize}
    \item \textbf{Pourcentages et coefficients multiplicateurs.}  
    Un taux d'évolution \(t\%\) se traduit par le coefficient multiplicateur
    \[
    CM = 1 + \frac{t}{100}.
    \]
    Si \(t>0\), \(CM>1\) : augmentation. Si \(t<0\), \(CM<1\) : diminution.
    \item \textbf{Relation fondamentale.}  
    \[
    V_f = CM \times V_i.
    \]
    On en déduit \(CM = \dfrac{V_f}{V_i}\) et \(V_i = \dfrac{V_f}{CM}\).
    \item \textbf{Taux d'évolution.}  
    \[
    t = \frac{V_f - V_i}{V_i} \times 100.
    \]
    \item \textbf{Évolutions successives.}  
    Les coefficients multiplicateurs se multiplient :
    \[
    CM_{\text{global}} = CM_1 \times CM_2 \times \cdots \times CM_n.
    \]
    Le taux global est \(T = (CM_{\text{global}} - 1) \times 100\).
    \item \textbf{Taux moyen.}  
    Pour \(n\) périodes, le taux moyen \(t_m\) vérifie
    \[
    \left(1 + \frac{t_m}{100}\right)^n = CM_{\text{global}}.
    \]
    Donc
    \[
    t_m = \left( (CM_{\text{global}})^{1/n} - 1 \right) \times 100.
    \]
    \item \textbf{Puissances et racines.}  
    \[
    a^{m/n} = \sqrt[n]{a^m}, \qquad (a^m)^n = a^{mn}, \qquad a^0 = 1.
    \]
    \item \textbf{Fractions.}  
    Diviser par un nombre revient à multiplier par son inverse :
    \[
    \frac{a}{b} = a \times \frac{1}{b} \quad (b \neq 0).
    \]
\end{itemize}

\subsection*{Exercice 1 -- Notions de base (4 points)}

\paragraph{1. Augmentation de \(15\%\).}
Le coefficient multiplicateur associé à une augmentation de \(15\%\) est
\[
CM = 1 + \frac{15}{100} = 1 + 0{,}15 = 1{,}15.
\]
Le nouveau prix est
\[
V_f = CM \times V_i = 1{,}15 \times 80 = 92.
\]
Le nouveau prix est donc \(92\)~\euro{}.

\paragraph{2. Taux d'évolution entre \(92\)~\euro{} et \(78{,}20\)~\euro{}.}
On a \(V_i = 92\) et \(V_f = 78{,}20\). Le taux d'évolution est
\[
t = \frac{V_f - V_i}{V_i} \times 100
= \frac{78{,}20 - 92}{92} \times 100
= \frac{-13{,}8}{92} \times 100.
\]
Or \(\frac{13{,}8}{92} = 0{,}15\). Donc
\[
t = -0{,}15 \times 100 = -15\%.
\]
Le signe est négatif : il s'agit d'une \textbf{baisse de \(15\%\)}.  
Vérification : \(92 \times (1 - 0{,}15) = 92 \times 0{,}85 = 78{,}20\).

\paragraph{3. Retrouver le prix initial après une baisse de \(8\%\).}
Une baisse de \(8\%\) correspond à
\[
CM = 1 - \frac{8}{100} = 1 - 0{,}08 = 0{,}92.
\]
On connaît \(V_f = 46\)~\euro{}. On cherche \(V_i\) tel que
\[
V_f = CM \times V_i.
\]
Donc
\[
V_i = \frac{V_f}{CM} = \frac{46}{0{,}92} = 50.
\]
Le prix initial était \(50\)~\euro{}.  
Vérification : \(50 \times 0{,}92 = 46\).

\subsection*{Exercice 2 -- Évolutions successives et taux global (6 points)}

\paragraph{1. Coefficients multiplicateurs.}
\begin{itemize}
    \item 2020 : perte de \(6\%\) donc \(CM_1 = 1 - 0{,}06 = 0{,}94\).
    \item 2021 : gain de \(10\%\) donc \(CM_2 = 1 + 0{,}10 = 1{,}10\).
    \item 2022 : perte de \(4\%\) donc \(CM_3 = 1 - 0{,}04 = 0{,}96\).
\end{itemize}

\paragraph{2. Coefficient multiplicateur global.}
On multiplie les coefficients :
\[
CM_{\text{global}} = CM_1 \times CM_2 \times CM_3
= 0{,}94 \times 1{,}10 \times 0{,}96.
\]
Calcul :
\[
0{,}94 \times 1{,}10 = 1{,}034,
\]
puis
\[
1{,}034 \times 0{,}96 = 0{,}99264.
\]
Donc \(CM_{\text{global}} = 0{,}99264\).

\paragraph{3. Taux global.}
\[
T = (CM_{\text{global}} - 1) \times 100
= (0{,}99264 - 1) \times 100
= -0{,}00736 \times 100
= -0{,}736\%.
\]
Le taux global est donc une \textbf{baisse de \(0{,}736\%\)}.

\paragraph{4. Nombre de clients fin 2022.}
\[
V_f = V_i \times CM_{\text{global}} = 15\,000 \times 0{,}99264 = 14\,889{,}6.
\]
En arrondissant à l'unité, on obtient \(14\,890\) clients.  
(On peut aussi faire les étapes successives :  
\(15\,000 \times 0{,}94 = 14\,100\),  
\(14\,100 \times 1{,}10 = 15\,510\),  
\(15\,510 \times 0{,}96 = 14\,889{,}6\).)

\paragraph{5. Pourquoi le stagiaire a tort.}
Le stagiaire additionne les taux : \(-6+10-4=0\). Mais les évolutions successives ne s'additionnent pas : elles se \textbf{composent} en multipliant les coefficients multiplicateurs.  
Ici :
\[
CM_{\text{global}} = 0{,}94 \times 1{,}10 \times 0{,}96 = 0{,}99264 \neq 1.
\]
Le nombre de clients n'est donc pas resté identique : il a légèrement baissé.

\subsection*{Exercice 3 -- Taux moyen et taux inconnu (10 points)}

\paragraph{1. Coefficient multiplicateur global.}
Les taux sont \(+10\%\), \(-5\%\), \(+8\%\), \(-3\%\), \(+6\%\).  
Les coefficients multiplicateurs correspondants sont :
\[
1{,}10,\quad 0{,}95,\quad 1{,}08,\quad 0{,}97,\quad 1{,}06.
\]
On les multiplie :
\[
CM_{\text{global}} = 1{,}10 \times 0{,}95 \times 1{,}08 \times 0{,}97 \times 1{,}06.
\]
Calcul pas à pas :
\[
1{,}10 \times 0{,}95 = 1{,}045,
\]
\[
1{,}045 \times 1{,}08 = 1{,}1286,
\]
\[
1{,}1286 \times 0{,}97 = 1{,}094742,
\]
\[
1{,}094742 \times 1{,}06 = 1{,}16042652.
\]
Donc
\[
CM_{\text{global}} \approx 1{,}16042652.
\]

\paragraph{2. Taux global.}
\[
T = (CM_{\text{global}} - 1) \times 100
= (1{,}16042652 - 1) \times 100
= 16{,}042652\%.
\]
Arrondi à \(0{,}01\%\), on obtient \(T \approx 16{,}04\%\).

\paragraph{3. Taux annuel moyen équivalent.}
Il y a \(n=5\) années. Le taux moyen \(t_m\) vérifie
\[
\left(1 + \frac{t_m}{100}\right)^5 = CM_{\text{global}}.
\]
Donc
\[
1 + \frac{t_m}{100} = (CM_{\text{global}})^{1/5}.
\]
On calcule
\[
(1{,}16042652)^{1/5} \approx 1{,}030205.
\]
Ainsi
\[
\frac{t_m}{100} = 0{,}030205
\quad\text{donc}\quad
t_m \approx 3{,}0205\%.
\]
Arrondi à \(0{,}01\%\), on trouve \(t_m \approx 3{,}02\%\).

\paragraph{4. Capital final.}
Le capital initial est \(V_i = 2\,000\)~\euro{}. Donc
\[
V_f = V_i \times CM_{\text{global}}
= 2\,000 \times 1{,}16042652
= 2\,320{,}85304.
\]
Arrondi au centime, le capital final est \(2\,320{,}85\)~\euro{}.

\paragraph{5. Taux annuel inconnu.}
Un placement augmente de \(44\%\) en deux ans avec le même taux annuel \(t\).  
Le coefficient multiplicateur global est
\[
CM_{\text{global}} = 1 + \frac{44}{100} = 1{,}44.
\]
Après deux années, on a
\[
\left(1 + \frac{t}{100}\right)^2 = 1{,}44.
\]
Donc
\[
1 + \frac{t}{100} = \sqrt{1{,}44} = 1{,}2.
\]
Ainsi
\[
\frac{t}{100} = 1{,}2 - 1 = 0{,}2
\quad\text{donc}\quad
t = 20\%.
\]
Le taux annuel est \(20\%\).

\end{document}